\scanpage{146}
\mspage{72}
\[
(240)\qquad\mathcal V_{D',Y}\unsure{\hookrightarrow|\longleftarrow}\underbrace{\check N_{D',Y}}_{\text{fibré normal de }D'\text{ dans }Y}
\]
\deleted{\unsure{transformant la}} dont le composé avec $Y\longrightarrow\mathcal V_{D',Y}$ soit la section nulle, et qui induise l'application identique sur les \deleted{\unsure{fib}}rés normaux le long de la section nulle. Ça \inserted{devrait} marche en tous cas dans le cas de vrais espaces analytiques paracompacts --- c'est faux par contre dans le cas de multiplicités (Ex. $\deleted{D'=\mathbb E^1_{\mathbb C}\subset Y=\mathbb E^2_{\mathbb C}}$, considérons $\mathbb E^1\subset\mathbb E^2$, avec l'opération du groupe discret $\mathbb Z$, \deleted{\unsure{dont}} le générateur opère par la transvection $\begin{pmatrix}1&\unsure{1|0}\\0&1\end{pmatrix}$, on pose $D'=(\mathbb E^1,\mathbb Z)$, $Y=(\mathbb E^2,\mathbb Z)$ --- la \deleted{fibré} extension des fibrés vectoriels différentiables définie par le plongement $D'\subset\mathcal V_{D',Y}\subset Y$ n'est pas scindable, alors que celle définie par la section nulle d'un fibré l'est...). Il est assez clair que lorsqu'on dispose d'une telle immersion (et sans condition $\unsure{=|C^{\infty}}$ de différentiabilité) celle-ci est une ``équivalence d'homotopie'', et
\margin{du moins pour des espaces analytiques ordinaires}
\margin{145}
